\documentclass[../Main.tex]{subfiles}

\begin{document}

\section{Schlussfolgerungen und weiteres Vorgehen} \label{sec:Schlussfolgerungen}

Damit sind wir am Schluss unserer Arbeit angekommen. In diesem Kapitel versuchen wir, unsere Erkenntnisfrage zu beantworten, und gehen dabei auf die Limitationen ein, die wir hatten. Zudem liefern wir weitere Ansätze zur Beantwortung unserer Erkenntnisfrage und schlagen Verbesserungen für unseren Prototyp vor.

\subsection{Beantwortung der Erkenntnisfrage} \label{sec:Beantwortung}

Unsere Erkenntnisfrage lautete:

\begin{quote}
    \textit{Wie müssen Gegnerverhalten und Gegnerplatzierung in einem 3D-Hack-and-Slash-Prototyp konzipiert sein, damit ein zusammenhängender Kampfabschnitt mit nachvollziehbarem Spielfluss, angemessener Schwierigkeit und befriedigendem Spielgefühl entsteht?}
\end{quote}

Im Folgenden fassen wir zusammen, was unsere acht Designvorschläge zur Beantwortung beitragen. Wir ordnen sie dabei danach, wie deutlich die Beobachtungen aus Abschnitt~\ref{sec:ErgVorschlag} für die Wirkung sprechen, die wir uns von ihnen erhofft hatten (Kapitel~\ref{sec:Umsetzung}).

Für das Gegnerverhalten zeigt sich eine starke Tendenz dafür, dass die Begrenzung der gleichzeitig erscheinenden Gegner (V2) die grossen Gegnermengen, die es in 3D-Hack-and-Slash-Spielen gibt, bekämpfbar macht und dass die Spieler durch die Kombination von Nah- und Fernkämpfern (V3) entscheiden müssen, welchen Gegner sie zuerst angreifen. Ebenso deutlich sprechen unsere Beobachtungen dafür, dass sich der Spieler als Grundzustand überlegen fühlt (V1) und dass das dynamische Verhalten der Gegner Druck erzeugt (V5).

Bei der Gegnerplatzierung zeigt sich eine starke Tendenz dafür, dass neue Gegner erkannt werden, wenn sie zuerst einzeln eingeführt werden, und dass ihre spätere Kombination die Schwierigkeit erhöht (V8). Eine ebenso starke Tendenz zeigt sich dafür, dass sich die Spieler in den Pausen zwischen den Räumen erholen (V7) und dass weite Räume mehr Bewegungsfreiheit und taktische Distanz bieten (V6).

Für das gezielte Durchbrechen der Überlegenheit (V1) zeigt sich dagegen eine erkennbare, aber keine starke Tendenz dafür, dass es Spannung erzeugt. Das liegt daran, dass das Durchbrechen nicht immer beabsichtigt war und nicht immer Spannung, sondern auch Frust erzeugt hat. Eine erkennbare Tendenz zeigt sich auch dafür, dass die Schwierigkeit taktisch statt mechanisch entsteht (V3), dass die Anfangsaufstellung Überblick gibt (V5), dass enge Räume zu methodischerem Spiel zwingen (V6), dass die Pausen zwischen den Wellen Entlastung bringen und die Schwierigkeit planbar ansteigt (V7) und dass die isolierte Einführung beim Verstehen neuer Gegner hilft (V8).

Bei der Ankündigung von Angriffen (V4) deutet sich an, dass bei den Nahkämpfern Animationen allein ausreichen, bei den Archers finden wir dagegen kaum Hinweise darauf. Ob eine Ankündigung über weitere Kanäle wie Sound bei den Archers ausreicht, konnten wir nicht überprüfen.

Durch all diese Designvorschläge zeigt sich eine starke Tendenz dafür, dass unsere Schwierigkeit angemessen war. Auch für einen gelungenen Spielfluss sprechen unsere Daten deutlich: Vier von fünf Testpersonen erlebten den Kampfabschnitt als zusammenhängend, und zwei spielten nach dem Ende sogar ungefragt weiter. Für ein befriedigendes Spielgefühl zeigt sich dagegen eine erkennbare, aber keine starke Tendenz, da es durch Faktoren beeinflusst wurde, die nicht direkt mit den Gegnern zu tun hatten, wie zum Beispiel, dass die Combos etwas zu träge waren (Abschnitt~\ref{sec:Gesamteindruck}). Auch der Spielfluss wurde an einer Stelle durch einen solchen Faktor gebrochen, weil das Ende eines Encounters nicht kommuniziert wurde.

Deshalb lautet die Antwort auf unsere Erkenntnisfrage vorläufig: Für unseren Prototyp deuten die Playtests darauf hin, dass Gegnerverhalten und Gegnerplatzierung vor allem dann zu einem nachvollziehbaren Spielfluss, einer angemessenen Schwierigkeit und einem befriedigenden Spielgefühl beitragen, wenn die Zahl der gleichzeitig aktiven Gegner begrenzt ist, Nah- und Fernkämpfer kombiniert werden, neue Gegner isoliert eingeführt und erst danach kombiniert werden und es zwischen den Räumen Pausen gibt. Die lesbare Anfangsaufstellung und die Pausen zwischen den Wellen konnten wir dagegen nur in Ansätzen belegen.

Diese Arbeit entstand aus dem Wunsch, Spiele nicht nur zu spielen, sondern sie genauer zu verstehen und selbst ein solches Spielerlebnis zu erschaffen. Das haben wir mit diesem Prototyp erreicht, und zumindest für unseren Kampfabschnitt konnten wir Tendenzen erkennen, wie Gegnerverhalten und Gegnerplatzierung zu einem solchen Spielerlebnis beitragen. Zudem betreffen diese Tendenzen speziell das 3D-Hack-and-Slash-Genre.

\subsection{Limitationen} \label{sec:Limitationen}

Im Ramen dieser Maturitätsarbeit unterlagen wir einigen limitationen, die die ergebnisse dieser arbeit einschränken. Im folgenden werden wir sie chronologisch auflisten:

für kapitel 3 die Referenzspielanalyse wurdne 2 spiele analysiert aber von zwei unterschiedlichen personen. die struktur des mda fragebogens ist zwar einheitlich aber einschätzungen wie befriedigen bleiben aber subjektiv. des weiteren hatten wir keinen zugang zu den quellcodes der spiele und dadurch keine tiefer gehenden einblicke wir konnten nur das beschreiben, was von aussen beobachtbar war.

für kapitel 4 die leveldesign recherche, beahndelt keines der quellen hack and slash genre direkt. alle erkenntnisse haben wir selbstständig auf unser gengre übertragen.

für kapitel 5, die umsetzung im prototyp. die exaten raumgrössen, schadenswerte und weitere parameter haben wir iterativ nach eigenem ermessen und spielgefühl festgelegt deshalb können wir nicht ausschliessen, dass diese auch in irgendeinder art einfluss auf unsere ergebnisse genommen haben. aufgrund von zeitmangel und priorisierung von systemen konnten wir niht alle designvorschläge 100 prozent wie vorgesehen umsetzen sondern nur teilweise. die erhaltenen resultate haben dennoch einen wert für weiteres vorgehen. die rede sit vor allem von der begrenzung der gleichzeitigen angreifer, ankündigung von angriffen über animation hinaus. faktoren ausserhalb der gegner haben auch unsere daten beeinflusst, trägheit der combos, dass türen aufgehen aber das nicht übermittelt wird und dass das zu verwirrung sorgt, das wenn man neu startet, dass man wieder am spawn ist.

kapitel 6, der playtest und die auswertung, mit fünf personen lassen sich keine representativen danten ermitteln deswegen konnten wir nnur tendenzen erkennen. und wo wir keine hinweise auf unsere teorien gefundne haben lässt sich uach nicht das gegneteil feststellen. die meisten testpersonen waren etwa in unserem alter also ende teenager anfang 20 und wir kannten sie alle um maximal 2 ecken., was die ehrlichkeit ihrer antworten beeinflusst haben könnte. dabei fanden zwei playtest über discord statt weshalb wir die mimik und körpersprache der personen nicht beobachten konnten. nach dem playtest sind uns zudem aufgefallen, das einige fragen die designvorscläge nicht genau messen beispielsweise ist q5 gemeint, das zwei frange gelichzeitig frägt. und weitere wie q6 q8 q10

des weiteren sind unsere ergebnisse nur abgesichert im bezug auf unseren eigenen prototyp und nicht generell alle hack-and-slash-spiele. besonder unsere intentionen für die drei qualitäten sind dabei ausschlaggebend. des weiteren haeb wir diesen prototyp selber entwickelt und deshalb könnten wir auch ienen bias bei der auswertung des playtest haben. wir haben auch nur ein set an playtests durchgeführt und in einer version, wir konnten die verbesserungsvorschläge aus 7.3 nicht umsetzen.

\textcolor{red}{Die Ergebnisse dieser Arbeit unterliegen mehreren Einschränkungen.}

\begin{itemize}[leftmargin=1.5cm]
    \item \textcolor{red}{[Stichpunkte für eure eigenen Absätze, sortiert nach Arbeitsschritt. Am Ende alles löschen, was ihr nicht braucht.]}
\end{itemize}

\textcolor{red}{\textbf{Referenzspielanalyse (Kapitel~\ref{sec:Referenzspielanalyse})}}
\begin{itemize}[leftmargin=1.5cm]
    \item \textcolor{red}{Jedes Spiel wurde nur von einer Person analysiert. Die Struktur des MDA-Fragebogens ist zwar einheitlich, Einschätzungen wie \enquote{befriedigend} bleiben aber subjektiv (Abschnitt~\ref{sec:Analysefragebogen}).}
    \item \textcolor{red}{Die Funktionsweise der Referenzspiele konnten wir nur aus dem beobachteten Verhalten erschliessen, da wir keinen Einblick in den Quellcode hatten (z.\,B.\ Rollensystem, Begrenzung der Angreifer in \game{NieR: Automata}).}
    \item \textcolor{red}{Pro Spiel haben wir nur einen Abschnitt analysiert, und diesen drei- bis fünfmal gespielt (Abschnitt~\ref{sec:Analysefragebogen}).}
\end{itemize}

\textcolor{red}{\textbf{Leveldesign-Recherche (Kapitel~\ref{sec:Leveldesign})}}
\begin{itemize}[leftmargin=1.5cm]
    \item \textcolor{red}{Keine der Quellen behandelt das Hack-and-Slash-Genre direkt. Wir haben sie selbst auf das Genre übertragen (Abschnitt~\ref{sec:VorgehenLeveldesign}).}
    \item \textcolor{red}{Es handelt sich überwiegend um Erfahrungswissen von Entwicklern, nicht um empirische Studien (Abschnitt~\ref{sec:VorgehenLeveldesign}).}
    \item \textcolor{red}{Wir haben nur wenige Quellen genutzt (zwei Fachbücher, drei GDC-Vorträge) und aus Zeitgründen nur ausgewählte Recherchepunkte in Designvorschläge übertragen.}
\end{itemize}

\textcolor{red}{\textbf{Umsetzung im Prototyp (Kapitel~\ref{sec:Umsetzung})}}
\begin{itemize}[leftmargin=1.5cm]
    \item \textcolor{red}{Raumgrössen, Schadenswerte und weitere Parameter haben wir iterativ nach unserem eigenen Spielgefühl festgelegt (Abschnitt~\ref{sec:VorgehenUmsetzung}).}
    \item \textcolor{red}{Einige Vorschläge sind nur teilweise umgesetzt und konnten deshalb wenig zeigen: Die Begrenzung der gleichzeitig angreifenden Gegner haben wir kaum genutzt (V2), Angriffe werden nur über Animationen angekündigt (V4), und die Gegner erscheinen erst nach dem Betreten eines Raumes (V5).}
    \item \textcolor{red}{Faktoren ausserhalb der Gegner haben die Ergebnisse beeinflusst: die Trägheit der Combos, die fehlende Rückmeldung beim Ende eines Encounters und an den Türen sowie der Neustart ganz am Anfang nach einem Tod (Abschnitt~\ref{sec:Gesamteindruck}).}
    \item \textcolor{red}{Der Kampfabschnitt ist linear, die Raumstruktur hatte deshalb kaum Einfluss auf den Spielfluss (Abschnitt~\ref{sec:Gesamteindruck}).}
    \item \textcolor{red}{Der Prototyp hat kein Tutorial. Die Aktionen haben wir den Testpersonen mündlich gezeigt, statt sie selbst entdecken zu lassen (Abschnitt~\ref{sec:Durchfuehrung}).}
\end{itemize}

\textcolor{red}{\textbf{Playtests und Auswertung (Kapitel~\ref{sec:Evaluation})}}
\begin{itemize}[leftmargin=1.5cm]
    \item \textcolor{red}{Mit fünf Testpersonen lassen sich nur Tendenzen erkennen (Abschnitt~\ref{sec:Auswertung}). Die Urteile in Kapitel~\ref{sec:Evaluation} sind deshalb als Tendenzen und nicht als Bestätigungen zu verstehen. Wo wir kaum oder keine Hinweise gefunden haben, bedeutet das nicht, dass das Gegenteil zutrifft. [Begriff: statt \enquote{kleine Stichprobe} besser \enquote{die geringe Zahl der Testpersonen}.]}
    \item \textcolor{red}{Die meisten Testpersonen kannten wir persönlich oder über Bekannte, was trotz unserer Bitte um ehrliche Antworten zu wohlwollenderen Einschätzungen geführt haben kann. Zudem sind alle in ihren späten Teenagerjahren oder Anfang zwanzig, und keine spielt ausschliesslich Hack-and-Slash-Spiele (Abschnitt~\ref{sec:AuswahlTP}).}
    \item \textcolor{red}{Zwei Playtests fanden über Discord statt, wo wir nur Bildschirm und Stimme wahrnehmen konnten, nicht aber Mimik und Körpersprache (Abschnitt~\ref{sec:Beobachtungsraster}).}
    \item \textcolor{red}{Alle Playtests hat dieselbe Person geleitet und beobachtet. Was notiert wurde, hängt deshalb auch von ihrer Aufmerksamkeit ab. In entspannten Phasen haben wir zudem mit den Testpersonen gesprochen (Abschnitt~\ref{sec:Durchfuehrung}).}
    \item \textcolor{red}{Die Protokolle sind nicht ganz vollständig: Einige Beobachtungen haben wir erst nachträglich aus der Erinnerung ergänzt, etwa den Tod von TP2 und dass TP3 und TP5 nach dem Ende weiterspielten. Da wir die Aufnahmen nach der Auswertung gelöscht haben, lässt sich das nicht mehr überprüfen. Bei TP4 waren einzelne Interviewfragen nur als Stichworte notiert.}
    \item \textcolor{red}{Einige Aussagen im Post-Questionnaire messen nicht genau, was der Vorschlag bewirken soll: Q5 fragt Anfangsaufstellung und Verhalten gleichzeitig ab (V5), Q6 fragt nur, ob sich enge und offene Räume anders anfühlten, nicht, ob sie schwerer waren (V6), Q8 fragt nur, ob die Staffelung bemerkt wurde, nicht, ob sie beim Verstehen half (V8), und Q10 fragt nach einem zusammenhängenden Spielfluss, nicht danach, ob man weiterspielen wollte (Spielfluss).}
    \item \textcolor{red}{Ohne Vergleichsversion lassen sich Effekte, die sich überlagern, kaum trennen, etwa ob der Druck vom dynamischen Verhalten (V5) oder vom gleichzeitigen Angreifen (V2) kam.}
\end{itemize}

\textcolor{red}{\textbf{Übergreifend (nicht einem Kapitel zuzuordnen)}}
\begin{itemize}[leftmargin=1.5cm]
    \item \textcolor{red}{Die Ergebnisse gelten nur für unseren Prototyp und einen einzelnen Kampfabschnitt. Ob sie sich auf andere Hack-and-Slash-Spiele übertragen lassen, können wir nicht sagen.}
    \item \textcolor{red}{Wir haben den Prototyp selbst entwickelt und die Playtests selbst ausgewertet. Unsere eigenen Erwartungen können Beobachtung und Interpretation beeinflusst haben.}
    \item \textcolor{red}{Im Rahmen der Maturitätsarbeit war nur ein Durchgang aus Umsetzung und Playtest möglich. Die Verbesserungen aus Abschnitt~\ref{sec:Verbesserungen} konnten wir nicht mehr testen.}
\end{itemize}

\subsection{Weiteres Vorgehen und Ausblick} \label{sec:Verbesserungen}

In einem weiteren Vorgehen würden wir weiterhin Gegnerverhalten und Gegnerplatzierung untersuchen. Zum einen würden wir einige unserer Designvorschläge, bei denen wir \textcolor{red}{nur eine erkennbare Tendenz oder kaum Hinweise} feststellen konnten, noch einmal und spezifischer überprüfen, damit wir diesmal klarere Antworten erhalten. \textcolor{red}{Zum anderen würden wir einzelne Vorschläge abwandeln. Für jeden Designvorschlag sieht das wie folgt aus.}

\paragraph{\textcolor{red}{V1: Überlegenheit als Grundzustand}}
Die Überlegenheit des Spielers konnten wir der Tendenz nach schon recht gut etablieren, sie aber noch nicht richtig geplant brechen \textcolor{red}{(Abschnitt~\ref{sec:ErgVorschlag}, V1)}. Die konkrete Verbesserung wäre dabei, den Heavy Soldier wieder etwas stärker zu machen. \textcolor{red}{Wir würden den Vorschlag weiterverfolgen, da unsere Daten zumindest in diese Richtung tendieren.}

\begin{itemize}[leftmargin=1.5cm]
    \item \textcolor{red}{[Aus V1 ergänzen:] Untersuchen, welche Arten von Brüchen der Überlegenheit Spass machen und welche nicht. Beispiel: Das Heilen hat als taktischer Bruch Spannung erzeugt, ob es Spass gemacht hat, ist offen. Die Trägheit der Combos war dagegen kein taktischer Bruch und eher frustrierend (Abschnitt~\ref{sec:ErgVorschlag}, V1).}
\end{itemize}

\paragraph{\textcolor{red}{V2: Begrenzung gleichzeitig aktiver Gegner}}
\textcolor{red}{Die Begrenzung der gleichzeitig angreifenden Gegner haben wir kaum genutzt. In einem nächsten Playtest würden wir sie gezielt einsetzen, auch um dem Stunlock entgegenzuwirken, der gerade dann entstand, wenn alle Gegner gleichzeitig angriffen (Abschnitt~\ref{sec:ErgVorschlag}, V2).}

\paragraph{\textcolor{red}{V3: Ein taktisches Dilemma durch Nah- und Fernkämpfer}}
\textcolor{red}{Das Dilemma wurde wahrgenommen, liess sich aber auch ignorieren (Abschnitt~\ref{sec:ErgVorschlag}, V3).}

\begin{itemize}[leftmargin=1.5cm]
    \item \textcolor{red}{[Fehlt:] Was würdet ihr bei V3 ändern oder genauer prüfen? Zum Beispiel: ob das Dilemma stärker wirkt, wenn Archers mehr Schaden machen und man sie deshalb nicht mehr ignorieren kann, oder eine Frage, die direkt nach taktischer statt mechanischer Schwierigkeit fragt.}
\end{itemize}

\paragraph{\textcolor{red}{V4: Klar telegrafierte Angriffe}}
Die Ankündigung von Angriffen würden wir mit mehreren Kanälen erneut testen\textcolor{red}{, etwa mit Sound und Partikeleffekten}. Dabei könnten wir vielleicht auch Unterschiede bei den Soldiers feststellen, ganz besonders aber bei den Archers. \textcolor{red}{Auch diesen Vorschlag würden wir weiterverfolgen, da unsere Daten zumindest bei den Nahkämpfern in diese Richtung tendieren.}

\paragraph{\textcolor{red}{V5: Lesbare Platzierung, dynamisches Verhalten}}
\textcolor{red}{In einem nächsten Schritt würden wir das dynamische Verhalten weglassen und nur die symmetrische Platzierung überprüfen.} Das ist uns dieses Mal nicht ganz gelungen, weil dieser Vorschlag zwei unterschiedliche Dinge zu überprüfen versuchte. Zudem könnte man prüfen, ob es einen Unterschied macht, ob die Gegner schon da sind, wenn man den Raum betritt, oder ob sie erst dann spawnen\textcolor{red}{, wie in \game{NieR: Automata} bzw.\ \game{Ninja Gaiden} (Abschnitt~\ref{sec:ErgVorschlag}, V5)}.

\paragraph{\textcolor{red}{V6: Der Raum als Druckmittel}}
\textcolor{red}{Die weiten Räume haben wie beabsichtigt gewirkt. Bei den engen Räumen ist die Wirkung dagegen weniger klar und hängt stark vom Spielstil ab. Den Vorschlag würden wir deshalb abgewandelt weiterverfolgen: Da der Druck in engen Räumen vor allem durch den fehlenden Platz zum Heilen entstand (Abschnitt~\ref{sec:ErgVorschlag}, V6), würden wir die Enge gezielter über die Zahl der Gegner steuern und nicht nur über die Raumgrösse. Zudem würden wir die Schwierigkeit der Räume direkt abfragen und nicht nur, ob sie sich anders anfühlten.}

\paragraph{\textcolor{red}{V7: Wellenstruktur mit Erholungsphasen}}
Auch die Wellenstruktur würden wir weiterverfolgen und dabei noch gezielter auf Ruhe und Anspannung achten, zum Beispiel mit einer Pulsmessung während des Playtests.\autocite[285--287]{drachen2018gamesuserresearch} Dabei würden wir zwischen Pausen unterscheiden, die wir dem Spieler vorgeben, und solchen, die sich der Spieler selbst verschafft. \textcolor{red}{Damit die Pausen zwischen den Wellen besser erkennbar sind, bräuchte es zudem eine Rückmeldung, wann eine Welle oder ein Encounter vorbei ist (siehe unten).}

\paragraph{\textcolor{red}{V8: Gestaffelte Einführung neuer Gegner}}
\textcolor{red}{Zwischen der Einführung eines neuen Gegners und seiner Kombination mit anderen würden wir mehr Zeit lassen, etwa wie in \game{NieR: Automata}, wo jede Vorstellungsphase ein ruhiges Tal vor der nächsten Arena bildet (Abschnitt~\ref{sec:ErgVorschlag}, V8). Zudem würden wir genauer abfragen, ob die isolierte Einführung beim Verstehen hilft, und nicht nur, ob sie bemerkt wurde.}

\paragraph{\textcolor{red}{Playtests}}
\textcolor{red}{Allgemein würden wir die Aussagen Q5, Q6, Q8 und Q10 beim nächsten Playtest genauer formulieren, da sie nicht immer genau das abfragten, was der jeweilige Vorschlag bewirken sollte (Abschnitt~\ref{sec:Limitationen}).} Ausserdem könnten wir Playtests mit einer grösseren Gruppe durchführen, um statistisch auswertbare Daten zu erhalten.

\paragraph{\textcolor{red}{Verbesserungen am Prototyp}}
Um dieses Vorgehen zu unterstützen \textcolor{red}{und Faktoren auszuschliessen, die nicht direkt mit den Gegnern zu tun haben}, würden wir an unserem Prototyp folgende Verbesserungen vornehmen:

\begin{itemize}[leftmargin=1.5cm]
    \item die Trägheit der Combos verringern (V1 und Abschnitt~\ref{sec:Gesamteindruck}),
    \item den Übergang zwischen Bewegung und Kampf verbessern (Abschnitt~\ref{sec:Gesamteindruck}),
    \item Audio- und/oder Partikeleffekte ergänzen, die Angriffe ankündigen (V4),
    \item einen Indikator einbauen, der anzeigt, wann ein Encounter vorbei ist, sowie Sounds für die Türen, damit auch diese klar erkennbar sind (Abschnitt~\ref{sec:Gesamteindruck}),
    \item I-Frames hinzufügen, etwa beim Aufstehen nach einem Treffer, um dem Stunlock entgegenzuwirken (V1 und V2),
    \item den Air Combat, der gegen unsere Bodengegner kaum nützlich war und trotzdem von einigen Testpersonen eingesetzt wurde (TP2, TP3 und TP4, jeweils B13), mit zwei Moves ergänzen: einer Aktion, die Gegner in die Luft wirft, und einer, mit der man aus der Luft auf den Boden dashen kann. Einen solchen Angriff aus der Luft wünschten sich auch mehrere Testpersonen (TP1, Q13; TP2, Q13; TP4, B11).
\end{itemize}

\end{document}
